\SourceLeaf{14}
signis superioribus vel inferioribus valentibus, prout $M$ iacet extra vel intra corpus. In casu priori autem partem soliditatis corporis constituunt partes coni a $P'$ usque ad $P''$, a $P'''$ usque ad $P''''$ etc., in posteriori vero partes coni a $M$ usque ad $P'$, a $P''$ usque ad $P'''$ etc. In utroque igitur casu pars corporis ea, quae iacet intra conum basi $\mathrm{d}\Sigma$ insistentem, fit
\[
=-\tfrac13(r'\mathrm{d}s'.\cos MQ'+r''\mathrm{d}s''.\cos MQ''+r'''\mathrm{d}s'''.\cos MQ'''+\text{ etc.})
\]
Tractando eodem modo cuncta elementa $\mathrm{d}\Sigma$, et summando, obtinemus
\par\medskip\Needspace{4\baselineskip}
\begin{center}THEOREMA QUINTUM.\end{center}
\par
\textit{Volumen corporis integri aequale est integrali $-\tfrac13\int r\mathrm{d}s.\cos MQ$ per totam corporis superficiem extenso.}

\Article{8}
\par
Iam supponamus, corpus esse uniformiter densum, singulaque eius elementa exercere attractionem in punctum $M$ alicui functioni distantiae proportionalem, ita ut denotante $\rho$ distantiam elementi a puncto attracto, attractio exprimatur per elementi volumen multiplicatum in $f\rho$. Concipiatur primo conus noster basi $\mathrm{d}\Sigma$ insistens totus materia plenus, atque per superficies sphaericas infinite sibi proximas e centro $M$ descriptas in elementa infinita dispertitus. Tale elementum, ad sphaeram cuius radius $=\rho$, exprimetur per $\rho\rho\mathrm{d}\rho.\mathrm{d}\Sigma$, adeoque vis, qua agit in $M$, per $\mathrm{d}\Sigma.\rho\rho f\rho.\mathrm{d}\rho$. Denotando itaque integrale $\int\rho\rho f\rho.\mathrm{d}\rho$ per $\Phi\rho$, patet $\mathrm{d}\Sigma.(\Phi\rho-\Phi0)$ exprimere attractionem partis coni a vertice usque ad distantiam $\rho$ in punctum $M$, sive generaliter $\mathrm{d}\Sigma.(\Phi\rho'-\Phi\rho)$ attractionem coni inter distantias a vertice $\rho$ et $\rho'$. Ab omnibus itaque partibus corporis nostri intra conum iacentibus attrahetur punctum $M$ in directione $M\Pi$ vi, quae exprimitur per
\[
\mathrm{d}\Sigma.(-\Phi r'+\Phi r''-\Phi r'''+\text{ etc.})
\]
quoties $M$ iacet extra corpus, vel per
\[
\mathrm{d}\Sigma.(-\Phi0+\Phi r'-\Phi r''+\Phi r'''-\text{ etc.})
\]
quoties $M$ iacet intra corpus, sive
\par
in casu priori per
